% !PRINTED 672
\weberpage{672}
%<<CONT
which we define by the equation

\weq{2}{\lambda_s = \delta_s \log|\eta_s|,}

where $|\eta_s|$ denotes the absolute value of $\eta_s$, and $\log|\eta_s|$ the real natural logarithm of the positive number $|\eta_s|$. If $\Omega_s$ is real and the sign $\pm$ is so chosen that $\pm \eta_s$ is positive, then

\[ \lambda_s = \log(\pm \eta_s). \]

But if $\eta_s, \eta_s'$ form an imaginary pair, then

\[ \lambda_s = \log \eta_s \eta_s', \]

and to $\eta_s$ and $\eta_s'$ there belongs the same $\lambda_s$, so that the number of distinct $\lambda_s$ is equal to $\nu$. From this determination there follows further

\weq{3}{\lambda_1 + \lambda_2 + \cdots + \lambda_\nu = \log N_a(\eta),}

where, as before, $N_a$ denotes the absolute norm.

The numbers $\lambda_1, \lambda_2, \ldots, \lambda_\nu$ we shall call the \gesperrt{conjugate logarithms} of the number $\eta$.

If $\eta$ is an \gesperrt{integer} of the field $\Omega$, then the absolute norm is a \gesperrt{natural number} which is $= 1$ only when $\eta$ is a unit, and from this it follows by (3):

1. \gesperrt{The sum of the conjugate logarithms of an integer $\eta$ is positive, and equal to zero only when $\eta$ is a unit.}

If $\omega_1, \omega_2, \ldots, \omega_n$ denotes a basis of the integers in $\Omega$ (a minimal basis of $\Omega$, \S.\,145), and $\omega_{1,s}, \omega_{2,s}, \ldots, \omega_{n,s}$ for $s = 1, 2, \ldots, n$ the conjugate numbers, then all integers of the field $\Omega_s$ can be represented in the form:

\weq{4}{\eta_s = x_1 \omega_{1,s} + x_2 \omega_{2,s} + \cdots + x_n \omega_{n,s},}

with integral rational coefficients $x_1, x_2, \ldots, x_n$. The determinant of this system

\[ \Sigma \pm \omega_{1,1} \omega_{2,2} \ldots \omega_{n,n} = \sqrt{\Delta} \]

is the square root of the discriminant $\Delta$ of the field, and therefore always different from zero. Accordingly the system (4) can be solved for the unknowns $x_1, x_2, \ldots, x_n$, and if we now stipulate that the absolute values of the numbers $\eta_s$ are not to exceed a certain limit, then from these solutions there result limits for the integral rational numbers $x_i$.
